\documentclass[10pt,a4paper]{amsart}
\usepackage[margin=19mm]{geometry}
\usepackage{amsmath,amssymb,mathtools}
\usepackage[scr=rsfs,cal=euler]{mathalfa}
\usepackage{xfrac}
\usepackage[unicode,hidelinks,bookmarks=false]{hyperref}
\newcommand{\ie}{i.e.\ }
\newcommand{\cf}{cf.\ }
\newcommand{\resp}{respectively\ }
\newcommand{\Subsection}{\subsection}
\providecommand{\nobreakdash}{\nobreak\hbox{-}}
\setlength{\emergencystretch}{2em}
\multlinegap0pt
\usepackage{amssymb}
\usepackage{algorithm}\usepackage{algpseudocode}
\newtheorem{theorem}{Theorem}
\newtheorem{lemma}{Lemma}
\newtheorem{observation}{Observation}
\title{Extremal structure in dense arrangements of $k$-intersecting curves}
\author{Andrew Suk, Su Zhou}
\date{Source: arXiv:2605.20705v1 (CC BY 4.0)}
\begin{document}
\maketitle

\noindent\emph{Source and licence.} Extremal structure in dense arrangements of $k$-intersecting curves. Version arXiv:2605.20705v1 (CC BY 4.0), distributed by arXiv under the Creative Commons Attribution 4.0 International licence, http://creativecommons.org/licenses/by/4.0/, as recorded in the arXiv licence metadata for this exact version and shown on the abstract page https://arxiv.org/abs/2605.20705v1. Full licence notice: \texttt{licenses/arxiv-2605.20705-CC-BY-4.0.txt}.\par\smallskip

\noindent\textit{Editorial scope.} Counted unit: the article's theorem soly (source0.tex lines 155--161). The article's complete original proof of it is delivered with it (source0.tex lines 998--1301, one \texttt{proof} environment of 304 lines, which the article places away from the statement). No mathematics is written, repaired or re-derived: the statement and the proof are the article's own lines, in the article's own notation, and no retained line is substituted or re-flowed.  Retained uncounted as the source-local context the proof needs: lines 295--312; lines 867--889; lines 902--915; lines 987--993.  One declared adaptation: exactly 1 ancillary strings inside the retained spans are omitted (image-inclusion commands and, where the source repeats a label, the repeated label definitions) and no other line differs from the source; the figure environments, with their placement, captions and labels, are kept, and the package records the omitted strings in fidelity receipts bound to the affected bodies.  No retained line contains a citation, so the wrapper carries no bibliography and no bibliography entry.  The retained lines contain the article's own diagram code, which the wrapper typesets with the standard tikz packages; no image file is delivered and the package declares no asset dependency.  Editorial formatting only, in addition: an AMS \texttt{amsart} preamble that restates the article's own declarations and macros used by the retained lines, namely the environments \texttt{theorem}, \texttt{lemma}, \texttt{observation} on separate counters, together with the standard packages those lines need.  The retained environments print in this document as Theorem 1, Lemma 1, Observation 1 in order of appearance, whereas the article prints the same items with the numbering of its own full document.  The complete native source of the article as distributed in the arXiv e-print for this exact version is bundled unchanged as \texttt{sources/201011/source0.tex}.
\begin{theorem}\label{rdiv}
There exists an absolute constant $r_0>0$ such that the following holds. Let $G$ be a planar graph on $N$ vertices embedded in the plane with all faces of size 2 or 3, and let $P \subseteq V(G)$.  Then for every $r \ge r_0$ and every $t \ge 1$, there exists a division $\mathcal R$ of $G$, obtained from a recursion tree $\mathcal T$ and a collection of simple cycles $\mathcal S$, such that
\[
|\mathcal{R}| = O\!\left(\frac{N}{r} + \frac{|P|}{t}\right),
\]
and every region $R \in \mathcal{R}$ satisfies $|V(R)| \le O(r),
|b(R)| = O(\sqrt r),$ and $|(V(R)\setminus b(R))\cap P| \le O(t).$  Moreover, every simple cycle $C_x \in \mathcal{S}$ satisfies one of the following: 


\begin{enumerate}
     
    \item the number of vertices of $G$ inside or on $C_x$ is at least $r/8$, or

    \item the number of boundary vertices inside or on $C_x$ is at least $\sqrt{r}/8$, or

    \item the number of points from $P$ inside or on $C_x$ is at least $t/8.$
 \end{enumerate}
\end{theorem}

\begin{lemma}[Pach--Sharir]\label{PSlemma}
Let \(P\) be a finite point set and $\mathcal C$ be a finite collection of $k$-intersecting curves.  Suppose there are $m$ crossing points between all pairs of curves in $\mathcal C$ outside of $P.$
Then
\[
I(P,\mathcal C)
\le
O_k\left(
|P|^{\frac{k + 1}{2k + 1}}
m^{\frac{k}{2k + 1}}
+|P|+|\mathcal  C|
\right).
\]
Since we have $m \leq k\binom{|\mathcal C|}{2}$, we obtain
\[
I(P,\mathcal{C}) = O_k\left(
|P|^{\frac{k + 1}{2k + 1}}
|\mathcal C|^{\frac{2k}{2k + 1}}
+|P|+|\mathcal  C|
\right).
\]  


\end{lemma}

\begin{lemma}\label{smalldeg}
Let $P$ be a set of points and $\mathcal C$ be a collection of $k$-intersecting curves in the plane. For $\ell \ge 2$,  we have


$$\sum\limits_{p\in P, |p| \geq \ell}|p| \leq  O_k\!\left(
\frac{|\mathcal C|^2}{\ell^{1+1/k}}
+
|\mathcal C|
\right).$$



 
\end{lemma}

\begin{lemma}\label{lem:supersat-hyper}
For every $k\ge 2$, $s\ge k+1$, and $\delta>0$ there exists
$\alpha=\alpha(k,s,\delta)>0$ such that the following holds.
If $\mathcal G$ is a $(k + 1)$-uniform hypergraph on $N$ vertices with at least
$\delta N^{k + 1}$ edge-disjoint copies of $K^{(k + 1)}_s$, then $\mathcal G$
contains at least $\alpha N^s$ copies of $K^{(k + 1)}_s$.
\end{lemma}

\begin{theorem}\label{soly}
Fix \(s>k+1\ge 2\), and let \(P\) be a set of \(n\) points in the plane and \(\mathcal C\) a set of \(n\) simple \(k\)-intersecting curves.
If there do not exist \(s\) points of \(P\) such that every \((k+1)\)-tuple among them is contained in a distinct curve of \(\mathcal C\), then
\[
I(P,\mathcal C)=o\!\left(n^{\frac{3k+1}{2k+1}}\right).
\]
\end{theorem}

\begin{proof}[Proof of Theorem~\ref{soly}]
Set $T=n^{\frac{1}{2k+1}}$. Let $P$ be a set of $n$ points and $\mathcal{C}$ be a collection of $n$ $k$-intersecting curves in the plane such that $I(P,\mathcal{C}) \geq \varepsilon n^{\frac{3k + 1}{2k + 1}}=\varepsilon \frac{n^2}{T^{k+1}}$, where $n > n_0,T>n^{\frac{1}{2k+1}}_0$ and $n_0$ is a sufficiently large constant that depends only on $\varepsilon, k, s$. We will show that, for sufficiently large \(n_0\), there exist \(s\) points of \(P\) such that every \((k + 1)\)-subset of them lies on a distinct curve of \(\mathcal C\). 


Without loss of generality, we can assume that every point in $P$ lies on at least one curve in $\mathcal{C}$, and each curve in $\mathcal{C}$ has at least one point from $P$ on it.  
Moreover, we can assume that the intersection graph of $\mathcal{C}$ is connected, since otherwise, we could pull apart the arrangement and work in each connected component individually.  

Set \(\ell=\frac{c_4 }{\varepsilon^{k/(k+1)}}T^k\), where \(c_4\) is a sufficiently large absolute constant to be chosen momentarily. Let \(Q\subset P\) be the set of points incident to at least \(\ell\) curves of \(\mathcal C\). Choosing \(c_4\) sufficiently large and applying Lemma~\ref{smalldeg}, we may ensure that the number of incidences between \(Q\) and \(\mathcal C\) is at most

$$I(Q,\mathcal C) \leq \sum\limits_{p \in Q}|p| \leq O\left(\frac{n^2}{\ell^{1 + 1/k}} + n\right) \leq \frac{\varepsilon}{2} \frac{n^2}{T^{k+1}}.$$  Thus, for $P'   = P\setminus Q$, the number of incidences between $P'$ and $\mathcal C$ is at least $\frac{\varepsilon}{2} \frac{n^2}{T^{k+1}}$.  Moreover, each point $p \in P'$ is incident to at most $\ell$ curves in $\mathcal{C}$.  Let us now restrict our attention to the arrangement $(P', \mathcal C)$.  By locally perturbing the curves in $\mathcal C$, without loss of generality, we can assume that no three curves in $\mathcal{C}$ have a common point outside of $P'$.
Let $m$ denote the number of crossing points among the curves in $\mathcal{C}$ outside of $P'$.  
By Lemma~\ref{PSlemma}, we have $m=\Omega_k\!\left(\varepsilon^{\frac{2k+1}{k}}n^2\right)$.  

From the arrangement $(P',\mathcal{C})$, consider the planar graph $G_0$ embedded in the plane whose vertex set consists of the $m$ crossing points together with the points of $P'$, and whose edges are consecutive vertices along a curve in $\mathcal{C}$.  
We then modify $G_0$ as follows.  For each point $p \in P'$, we add $w = \varepsilon^{(2k + 1)/k}T^{2k+1}/\ell=\varepsilon^{(2k + 1)/k}n/\ell$ nested cycles centered at $p$, and very close to $p$, such that each such cycle has length $2|p|$ and the resulting graph remains planar.   See Figure~\ref{X}.  





 

\begin{figure}
    \centering
    
    \caption{The local modification near $p$.}
    \label{X}
\end{figure}

After performing this operation for every $p \in P'$, and since $|p|\leq \ell$,  we will have added at most an additional $2\varepsilon^{\frac{2k+1}{k}}n^2$ vertices to $G_0$. We then triangulate every face of size greater than three in the resulting graph, obtaining a planar graph \(G=(V,E,F)\) embedded in the plane in which every face has size \(2\) or \(3\), and
\[
|V| = m + 2\varepsilon^{\frac{2k+1}{k}}n^2 + |P'| = O(m).
\]

\noindent  Note that $ m \leq kn^2$.  We now apply Theorem~\ref{rdiv} to $G$ and the  point set $P'\subset V(G)$, with parameters











\[
r = c_5\frac{k^2s^2}{\varepsilon^2}T^{2k+2}, 
\qquad
t =  c_5\frac{ks^2}{\varepsilon^2}T,
\]

\noindent where $c_5$ is a large absolute constant that will be chosen later.  Hence, we obtain a division $\mathcal R$ such that

$$|\mathcal R| = O\left( \frac{|V|}{r } + \frac{n}{t}\right) =  
O\!\left(
\frac{\varepsilon^2}{c_5ks^2}
\frac{n^2}{T^{2k+2}}
\right).$$


\noindent and for each region $R_i \in \mathcal R$ we have


\[
|V(R)|\le r,\qquad |b(R)|=O(\sqrt r), \qquad 
  |(V(R)\setminus b(R))\cap P'| \le t.
\]

\noindent  Moreover, there is a collection of simple cycles (simple closed curves in the plane) $\mathcal S$ such that for each boundary vertex $v$, there is a simple cycle $C_x \in \mathcal{S}$ that contains $v$, such that either


\begin{enumerate}
     
    \item the number of vertices of $G$ inside or on $C_x$ is at least $r/8$, or

    \item the number of boundary vertices inside or on $C_x$ is at least $\sqrt{r}/8$, or

    \item the number of points from $P$ inside or on $C_x$ is at least $t/8.$
 \end{enumerate}

 

Let \(Q' \subset P'\) be the set of boundary vertices of the division \(\mathcal R\).  That is, \(p\in Q'\) implies that \(p\) lies on some cycle \(C_i\in\mathcal S\). Fix \(p\in Q'\), and let \(Z_1,\dots,Z_w\) be the $w$ nested cycles constructed around $p$.  The cycles are ordered from inside to outside, so that \(Z_1\) is the innermost and \(Z_w\) the outermost. Let \(G_p\) be the subgraph of \(G\) induced by \(p\) together with the vertices of \(Z_1,\dots,Z_w\).  See Figure \ref{X}.














\begin{observation}
If $p \in Q'$, then at least  $c_6T^{k + 1}$ vertices from $G_p$ are also boundary vertices, where $c_6 = c_6(k,\varepsilon)$.
\end{observation}





\begin{proof}
 
The proof falls into two cases.

\medskip
\noindent\emph{Case 1.} Suppose that there exists a node \(x\in\mathcal T\) such that the separating cycle $C_x$ lies inside of $Z_i$ for some $i\geq 1$. That is, the simple closed curve in the plane corresponding to $C_x$ lies inside the simple closed curve corresponding to $Z_i$.  Hence, $V(C_x) \subset V(G_p)$.   Since \(G_p\) has \(1+2w|p|\) vertices and \(|p|\le \ell\), we have
\[
|V(G_p)| \le 1+2w\ell = O\!\left(\varepsilon^{(2k+1)/k}T^{2k+1}\right).
\]
For \(T\) sufficiently large,
\[
|V(G_p)|<\frac r8,
\qquad\text{and}\qquad
|V(G_p)\cap P'|=1<\frac t8.
\]
Therefore, property (2) in Theorem \ref{rdiv} must hold, and the number of boundary vertices inside or on \(C_x\) is at least $\frac{\sqrt r}{8}
=
\frac{\sqrt{c_5}\,ks}{8\varepsilon}\,
T^{k+1}.$  Since \(C_x\) lies inside $Z_i$, all of these boundary vertices lie in \(G_p\). Thus for $c_6 > 0$ sufficiently small, \(G_p\) contains at least $c_6 T^{k+1}$ boundary vertices.

\medskip
\noindent\emph{Case 2.} Suppose we are not in Case 1.  Let \(x\in\mathcal T\) be the highest node such that \(p\in V(C_x)\), and let \(R_x\) be the corresponding region. 

Let \(i\) be the largest index such that the cycles \(Z_1,\dots,Z_i\) are fully contained in \(R_x\). (Possibly \(i=0\).) Then for each \(j>i\), the cycle \(Z_j\) is not fully contained in \(R_x\). We first claim that for every \(j>i\), the cycle \(Z_j\) already contains a boundary vertex. Indeed, \(Z_j\) is fully present in the root region, but not fully present in \(R_x\). Hence along the root-to-\(x\) path in \(\mathcal T\), there is a first node \(y_j\) such that \(Z_j\) is not fully contained in \(R_{y_j}\). Let \(z_j\) be the parent of \(y_j\). Then \(Z_j\) is fully contained in \(R_{z_j}\), but not in \(R_{y_j}\). Therefore the separator cycle \(C_{z_j}\) must contain a vertex of \(Z_j\). Such a vertex is a boundary vertex. Since the cycles \(Z_j\) are pairwise disjoint, these boundary vertices are distinct for different \(j>i\). Thus we obtain at least \(w-i\) boundary vertices from the cycles \(Z_{i+1},\dots,Z_w\).

Next consider the cycles \(Z_1,\dots,Z_i\). Since each such cycle is fully contained in \(R_x\), while \(C_x\) contains a point outside of $Z_i$, the Jordan curve \(C_x\) must meet each \(Z_1,\ldots, Z_i\). Hence for each \(1\le j\le i\), the cycle \(Z_j\) contains a boundary vertex lying on \(C_x\). Again these vertices are distinct because the cycles \(Z_1,\dots,Z_i\) are pairwise disjoint.

Combining the two parts, we find that \(G_p\) contains at least
\[
(w-i)+i=w
\]
distinct boundary vertices. Since
\[
w=\frac{\varepsilon^{(2k+1)/k}T^{2k+1}}{\ell}
\geq c_6 T^{k+1},
\]
by setting \(c_6=c_6(k,\varepsilon)\) sufficiently small, the proof is complete.
\end{proof}



























By the Observation above, the total number of boundary vertices is at least $|Q'|c_6T^{k + 1}$, which implies


$$|Q'|c_6T^{k + 1} \leq O( |\mathcal R|\sqrt{r}) \leq O\!\left(
\frac{\varepsilon}{\sqrt{c_5}\,s}
\frac{n^2}{T^{k+1}}
\right).
$$


\noindent  This implies

\[
|Q'|
\le
O\!\left(
\frac{\varepsilon}{c_6\sqrt{c_5}\,s}
\frac{n^2}{T^{2k+2}}
\right)=O\!\left(
\frac{\varepsilon}{c_6\sqrt{c_5}\,s}
\frac{n}{T}
\right).
\]

\noindent By Lemma \ref{PSlemma}, we have

$$ I(Q',\mathcal C) = O_k\left(
|Q'|^{\frac{k + 1}{2k + 1}}
|\mathcal C|^{\frac{2k}{2k + 1}}
+|Q'|+|\mathcal  C|
\right) \leq  O_k\!\left(
\left(\frac{\varepsilon}{c_6\sqrt{c_5}\,s}\right)^{\frac{k+1}{2k+1}}
\frac{n^2}{T^{k+1}}\frac{1}{T^{\frac{k+1}{2k+1}}}
+
\frac{\varepsilon }{c_6\sqrt{c_5}\,s}\frac{n^2}{T^{2k+2}}
+
n
\right).$$

\noindent Hence for $n_0 = n_0(c_5,c_6,s, \varepsilon, k)$ sufficiently large, for all $n > n_0$ with $T>n^{\frac{1}{2k+1}}_0$ we have



\[
I(Q',\mathcal C) \leq \frac{\varepsilon}{4}\frac{n^2}{T^{k+1}}.\]



Hence, there are at least $\frac{\varepsilon}{4}\frac{n^2}{T^{k+1}}$ incidences between $P'\setminus Q'$ and $\mathcal C$.  The division $\mathcal R$ gives a natural partition of $P'\setminus Q' = \bigcup_i P_i$, where 

$$P_i = \{p \in P'\setminus Q': p \in V(R_i)\}.$$

\noindent  Note that all points in $P_i$ are not boundary vertices.   By Theorem~\ref{rdiv}, we have that $|P_i|  \leq t =  c_5\frac{ks^2}{\varepsilon^2}T$. 




Let $\Gamma$ be the incidence graph between $P'\setminus Q'$ and $\mathcal C$. Consider a curve $\gamma \in \mathcal{C}$.  Let $b(\gamma)$ be the number of boundary vertices of $G$ on $\gamma$ with respect to $\mathcal{R}$.  Along each maximal subarc of \(\gamma\) whose interior contains no boundary vertex, the points from a fixed part \(P_i\) appear consecutively, we partition them into consecutive blocks of size \(s\), and discard the remaining incidences from the subarc which has at most $s-1$ points. Since each such subarc lies between two consecutive boundary points on \(\gamma\), the total number of deleted incidences on \(\gamma\) is at most \(s\,(b(\gamma)+1)\).  Each boundary vertex either is a point in $Q'$ with degree at most $2\ell$, or is a vertex with degree 4. Hence, after performing this operation for each $\gamma \in \mathcal C$, we will delete at most
\[
\sum\limits_{\gamma \in \mathcal C} s(b(\gamma)  +1)  = O\!\left(s\sqrt r\,|\mathcal R|+s\ell|Q'|+s|\mathcal C|\right)
=
O\!\left(
\frac{\varepsilon}{\sqrt{c_5}}
\frac{n^2}{T^{k+1}}
+
\frac{\varepsilon^{\frac{1}{k+1}}}{\sqrt{c_5}}
\frac{n^2}{T^{k+2}}+n
\right).
\]
 
\noindent incidences.  Indeed, there are at most $O(\sqrt{r}|\mathcal R|)$ boundary points, and each point $p \in Q'$ has at most $\ell$ curves going through it.   By setting $c_5$ to be a sufficiently large absolute constant, and for $n,T$ sufficiently large, the quantity above is less than $\frac{\varepsilon}{8} \frac{n^2}{T^{k+1}}$. Thus, at least $\frac{\varepsilon}{8} \frac{n^2}{T^{k+1}}$ edges (incidences) remain in $\Gamma$. 

By the pigeonhole principle, there is a part $P_i \subset P'\setminus Q'$ such that the number of edges between $P_i$ and $\mathcal C$ in $\Gamma$ is at least
\[
\frac{\varepsilon \frac{n^2}{T^{k+1}}}{8|\mathcal R|}
=
\Omega\!\left(
\frac{ks^2}{\varepsilon}
T^{k+1}
\right),
\]

\noindent and
\[
 |P_i|  \leq t \leq  c_5\frac{ks^2}{\varepsilon^2}T.
\]

Moreover, if a curve $\gamma \in \mathcal C$ is incident to a point $p \in P_i$, and $(p,\gamma)$ is a surviving edge in $\Gamma$, then $\gamma$ must be incident to at least $s$ vertices in $P_i$.  Let $\mathcal C_0$ be the curves in $\mathcal C$ that are incident to at least $s$ points in $P_i$. By Lemma~\ref{PSlemma}, we have
\[
\Omega\!\left(
\frac{ks^2}{\varepsilon}
T^{k+1}
\right)  = O_k\left(
|P_i|^{\frac{k + 1}{2k + 1}}
|\mathcal C_0|^{\frac{2k}{2k + 1}}
+|P_i|+|\mathcal  C_0|
\right).
\]

\noindent Since $|P_i|=O_{\varepsilon,k}\!\left(T\right)$, the term $|P_i|$ is negligible for large $T$. Therefore 


\[
|\mathcal C_0|
=
\Omega_k\!\left(
 \sqrt{k}\,s\,\varepsilon^{1/(2k)} 
T^{k+1}
\right).
\]

Let $H=(P_i,E)$ be the auxiliary $(k + 1)$-uniform hypergraph whose vertex set is $P_i$. For each $\gamma \in \mathcal C_0$, consider the points of $P_i$ that appear along $\gamma$ from one end to the other. We then create a $K^{(k+1)}_s$ on the first $s$ points along $\gamma$, then create another $K^{(k+1)}_s$ on the next set of $s$ points along $\gamma$, and repeat this process until there are fewer than $s$ points left on $\gamma$. Thus each $\gamma$ gives rise to at least one copy of $K^{(k+1)}_s$.  Moreover, since no two curves in $\mathcal C_0$ have $k + 1$ points in common, the corresponding copies of $K_s^{(k + 1)}$ are edge disjoint. Since
\[
|\mathcal C_0| \geq \Omega_k\!\left(
 \sqrt{k}\,s\,\varepsilon^{1/(2k)} 
T^{k+1}
\right)
\quad\text{and}\quad
N = |P_i|  \leq c_5\frac{ks^2}{\varepsilon^2}T,
\]
we have $|\mathcal C_0| \geq \delta N^{k + 1}$ for some constant $\delta = \delta(\varepsilon,k,s) > 0$. Thus $H$ is a $(k +1)$-uniform hypergraph on $N$ vertices that contains at least $\delta N^{k + 1}$ edge-disjoint copies of $K^{(k+1)}_s$. By Lemma~\ref{lem:supersat-hyper}, it follows that $H$ contains at least $\alpha N^s$ copies of $K^{(k+1)}_s$, where $\alpha=\alpha(k,\varepsilon,s)$.

We now bound the number of $K^{(k+1)}_s$ whose $s$ vertices contain $k + 2$ points from $P_i$ lying on a single curve $\gamma\in\mathcal C_0$. These $k + 2$ points must lie in a single block of $s$ consecutive points along that curve.  There are at most $sN^{k + 1}$ such $(k +2)$-tuples, since no two curves have $k + 1$ points in common. Once the \(k+2\) vertices lying on a single curve are chosen, the remaining \(s-k-2\) vertices may be chosen arbitrarily from \(P_i\). Thus each such $(k+2)$-tuple gives rise to at most \(O(N^{\,s-k-2})\) bad copies. Since there are at most \(N^{k+1}\) such tuples, the total number of such copies of $K_s^{(k + 1)}$ is
\[
O\!\left(sN^{k+1}\cdot N^{\,s-k-2}\right)=O(sN^{s-1}).
\]

 Finally, by setting $n > n_0$ sufficiently large, the number of such bad copies is at most
\[
csN^{s-1} < \alpha N^s,
\]
since \(N\) will be large. Hence, there exists a copy of \(K^{(k + 1)}_s\) in $H$ in which no $k + 2$ vertices lie on the same curve of \(\mathcal C_0\). Equivalently, there are \(s\) points of \(P_i\) such that each $(k + 1)$-tuple  is joined by a distinct curve from \(\mathcal C_0\), yielding \(\binom{s}{k + 1}\) distinct curves determined by these points.
\end{proof}

\end{document}
